% Standalone paragraph. Paste into your own thesis, or \input it.
% Not wired into any chapter's numbering -- drop it wherever the "whole tree hasn't been viewed
% yet" argument is needed.

% HEADLINE: one of three sibling subsections for the State of the Art chapter, sitting under
% \section{The gap} (state_of_the_art/08_the_gap.tex): where the tree is measured
% (whole_tree_gap), how it is measured (reproducible_features), in whom it is measured
% (healthy_cohort_2). Promote to \section if they are to stand as sections of their own.
\subsection{Where the tree is measured}
\label{sec:gap-where}

Whole coronary trees are rarely the unit of analysis.
The largest whole-tree geometry study identified in this literature covers 39 patients, and its
own authors frame that scope as the exception rather than the norm \cite{zhang2024curvature}.

That concentration is not accidental.
The left main is short and proximal, and an occlusion there cuts off both the LAD and circumflex
territories at once --- more myocardium than any single branch further out could threaten on its
own. Guidelines reflect that stake directly: a left main stenosis is classified as severe at 50\%,
where every other coronary segment needs 70\% to earn the same label \cite{knuuti2019esc}. Given
that, decades of feature-extraction work concentrating on one junction is not a mystery.

But the field's own recent review names the consequence of that choice as one of its central
limitations. A 2026 review of coronary anatomy and haemodynamics states plainly that most studies
covered the main bifurcation or non-bifurcating vessels, leaving most curved segments unmeasured,
and calls investigating the whole tree essential \cite{shen2026geometry}. What that leaves out is
exactly where two well-documented risks sit: the plaque phenotype most likely to rupture clusters
in the proximal left anterior descending artery, with further peaks in the right coronary and
circumflex arteries \cite{araki2020plaques}, and spontaneous coronary artery dissection occurs
predominantly in the middle-to-distal segments of the tree \cite{inoue2021scad}. Neither sits at
the left main --- the one segment most studies actually measure.

Analysing the whole tree, rather than inferring its condition from one focal measurement, is
therefore likely to reveal more than the left main can show on its own.
